% ==================================================================
% id:       app-022
% topic:    שימושי הנגזרת
% subtopic: חקירת פונקציה
% solution: full
% verified: no
% ==================================================================
\textbf{חקרו את הפונקציות הבאות (מצאו תחום הגדרה, בדקו זוגיות, אסימפטוטות, נקודות חיתוך עם
הצירים, תחומי עלייה וירידה ונקודות קיצון, ציירו סקיצה של גרף):}
\bE\bMs{2}
\item $\ds{f(x) = \frac{2 - x^2}{1 + x^4}} $
\hint{בדקו קודם את הזוגיות: לפונקציה זוגית מספיק לחקור את $x\ge0$. כדי למצוא את אפסי הנגזרת הציבו $t=x^2$.}
\begin{solbox}
\unverified
\textbf{תחום הגדרה:}
המכנה חיובי תמיד, ולכן
$D=\rr$.

\textbf{זוגיות:}
$f(-x)=\frac{2-x^2}{1+x^4}=f(x)$,
ולכן
$f$
זוגית (הגרף סימטרי ביחס לציר
$y$).

\textbf{חיתוך עם הצירים:}
$f(0)=2$,
כלומר
$(0,2)$.
$f(x)=0$
אמ"מ
$x^2=2$,
כלומר
$(\pm\sqrt2,0)$.

\textbf{אסימפטוטות:}
אין אסימפטוטות אנכיות כי
$f$
רציפה על כל הישר. מתקיים
$$\lim_{x\to\pm\infty}\frac{2-x^2}{1+x^4}=\lim_{x\to\pm\infty}\frac{\frac{2}{x^4}-\frac{1}{x^2}}{\frac{1}{x^4}+1}=0$$
ולכן
$y=0$
אסימפטוטה אופקית ב-$\pm\infty$.

\textbf{עלייה, ירידה וקיצון:}
\begin{align*}
f'(x)&=\frac{-2x(1+x^4)-(2-x^2)\cdot 4x^3}{(1+x^4)^2}
=\frac{2x^5-8x^3-2x}{(1+x^4)^2}
=\frac{2x\lr{x^4-4x^2-1}}{(1+x^4)^2}.
\end{align*}
הביטוי
$x^4-4x^2-1$
מתאפס כאשר
$x^2=2+\sqrt5$
(הפתרון
$x^2=2-\sqrt5<0$
נפסל), כלומר עבור
$x=\pm a$
כאשר
$a=\sqrt{2+\sqrt5}\approx 2.06$.
הוא שלילי עבור
$|x|<a$
וחיובי עבור
$|x|>a$.
לכן:
\begin{itemize}
\item
$f'<0$
ב-$(-\infty,-a)$
וב-$(0,a)$
-- שם
$f$
יורדת.
\item
$f'>0$
ב-$(-a,0)$
וב-$(a,\infty)$
-- שם
$f$
עולה.
\end{itemize}
מכאן ש-$x=0$
נקודת מקסימום מקומי (וגם מוחלט), עם
$f(0)=2$,
ו-$x=\pm a$
נקודות מינימום מקומי (וגם מוחלט). מכיוון ש-$a^4=(2+\sqrt5)^2=9+4\sqrt5$,
$$.f(\pm a)=\frac{2-(2+\sqrt5)}{10+4\sqrt5}=\frac{-\sqrt5}{10+4\sqrt5}=1-\frac{\sqrt5}{2}\approx-0.12$$

\textbf{סקיצה:}
הגרף מגיע מ-$0$ (מתחת לציר) ב-$-\infty$, יורד למינימום ב-$x=-a$, עולה וחותך את ציר
$x$
ב-$-\sqrt2$
עד למקסימום
$(0,2)$,
ובאופן סימטרי בצד ימין.
\end{solbox}
\item $\ds{f(x) = (x - 3)\sqrt{x}}$
\begin{solbox}
\unverified
\textbf{תחום הגדרה:}
$D=[0,\infty)$.

\textbf{זוגיות:}
התחום אינו סימטרי, ולכן
$f$
אינה זוגית ואינה אי-זוגית.

\textbf{חיתוך עם הצירים:}
$f(0)=0$,
ו-$f(x)=0$
אמ"מ
$x=0$
או
$x=3$.
כלומר
$(0,0)$
ו-$(3,0)$.

\textbf{אסימפטוטות:}
אין אסימפטוטות אנכיות
($f$
רציפה בתחום הגדרתה). כאשר
$x\to\infty$
מתקיים
$f(x)\to\infty$
וגם
$\frac{f(x)}{x}=\frac{x-3}{\sqrt x}\to\infty$,
ולכן אין אסימפטוטה אופקית או משופעת.

\textbf{עלייה, ירידה וקיצון:}
עבור
$x>0$
$$.f'(x)=\sqrt{x}+\frac{x-3}{2\sqrt{x}}=\frac{2x+x-3}{2\sqrt{x}}=\frac{3(x-1)}{2\sqrt{x}}$$
לכן
$f'<0$
ב-$(0,1)$
ו-$f'>0$
ב-$(1,\infty)$:
$f$
יורדת ב-$[0,1]$
ועולה ב-$[1,\infty)$.
הנקודה
$x=1$
היא מינימום מקומי (וגם מוחלט) עם
$f(1)=-2$,
ו-$x=0$
היא נקודת קצה שבה מתקבל מקסימום מקומי (חד-צדדי)
$f(0)=0$.
נשים לב כי
$f'(x)\to-\infty$
כאשר
$x\to0^+$,
כלומר הגרף יוצא מהראשית עם משיק אנכי.

\textbf{סקיצה:}
הגרף יוצא מהראשית כלפי מטה, יורד עד
$(1,-2)$,
עולה, חותך את ציר
$x$
ב-$(3,0)$
ושואף ל-$\infty$.
\end{solbox}
\item $\ds{f(x) = \frac{\ln x}{\sqrt{x}}}$
\begin{solbox}
\unverified
\textbf{תחום הגדרה:}
$D=(0,\infty)$.

\textbf{זוגיות:}
התחום אינו סימטרי, ולכן
$f$
אינה זוגית ואינה אי-זוגית.

\textbf{חיתוך עם הצירים:}
אין חיתוך עם ציר
$y$
($0\notin D$).
$f(x)=0$
אמ"מ
$x=1$,
כלומר
$(1,0)$.

\textbf{אסימפטוטות:}
כאשר
$x\to0^+$
מתקיים
$\ln x\to-\infty$
ו-$\frac{1}{\sqrt x}\to\infty$,
ולכן
$f(x)\to-\infty$:
הישר
$x=0$
אסימפטוטה אנכית. כאשר
$x\to\infty$,
לפי כלל לופיטל,
$$\lim_{x\to\infty}\frac{\ln x}{\sqrt{x}}=\lim_{x\to\infty}\frac{\frac1x}{\frac{1}{2\sqrt x}}=\lim_{x\to\infty}\frac{2}{\sqrt x}=0$$
ולכן
$y=0$
אסימפטוטה אופקית ב-$\infty$.

\textbf{עלייה, ירידה וקיצון:}
$$.f'(x)=\frac{\frac{1}{x}\sqrt{x}-\ln x\cdot\frac{1}{2\sqrt{x}}}{x}=\frac{2-\ln x}{2x\sqrt{x}}$$
המכנה חיובי, ולכן
$f'>0$
אמ"מ
$\ln x<2$,
כלומר
$x<e^2$.
לכן
$f$
עולה ב-$(0,e^2]$
ויורדת ב-$[e^2,\infty)$,
ו-$x=e^2$
נקודת מקסימום מקומי (וגם מוחלט) עם
$$.f(e^2)=\frac{2}{e}$$

\textbf{סקיצה:}
הגרף עולה מ-$-\infty$ (ליד ציר
$y$),
חותך את ציר
$x$
ב-$(1,0)$,
מגיע למקסימום
$\lr{e^2,\frac2e}$
ויורד בחזרה אל
$0$
מלמעלה.
\end{solbox}
\item $\ds{f(x) = (x + 2)e^{\frac{1}{x}}}$
\hint{בדקו בנפרד את הגבולות החד-צדדיים ב-$0$. בחיפוש אסימפטוטה משופעת היעזרו בגבול $\lim_{t\to0}\frac{e^t-1}{t}$.}
\begin{solbox}
\unverified
\textbf{תחום הגדרה:}
$D=\rr\setminus\{0\}$.

\textbf{זוגיות:}
$f(1)=3e$
ו-$f(-1)=e^{-1}$,
ולכן
$f$
אינה זוגית ואינה אי-זוגית.

\textbf{חיתוך עם הצירים:}
אין חיתוך עם ציר
$y$.
מכיוון ש-$e^{\frac1x}>0$,
$f(x)=0$
אמ"מ
$x=-2$,
כלומר
$(-2,0)$.

\textbf{אסימפטוטות:}
\begin{itemize}
\item
כאשר
$x\to0^+$
מתקיים
$\frac1x\to\infty$,
ולכן
$f(x)\to 2\cdot\infty=\infty$:
הישר
$x=0$
אסימפטוטה אנכית (מימין).
כאשר
$x\to0^-$
מתקיים
$e^{\frac1x}\to0$,
ולכן
$f(x)\to0$.
\item
אסימפטוטה משופעת
$y=mx+n$
ב-$\pm\infty$:
$$,m=\lim_{x\to\pm\infty}\frac{f(x)}{x}=\lim_{x\to\pm\infty}\lr{1+\frac2x}e^{\frac1x}=1$$
$$.n=\lim_{x\to\pm\infty}\lr{f(x)-x}=\lim_{x\to\pm\infty}\lr{x\lr{e^{\frac1x}-1}+2e^{\frac1x}}=1+2=3$$
כאן השתמשנו ב-$\lim_{t\to0}\frac{e^t-1}{t}=1$
עם
$t=\frac1x$.
לכן
$y=x+3$
אסימפטוטה משופעת ב-$\infty$
וב-$-\infty$.
\end{itemize}

\textbf{עלייה, ירידה וקיצון:}
$$.f'(x)=e^{\frac1x}+(x+2)e^{\frac1x}\cdot\lr{-\frac{1}{x^2}}=e^{\frac1x}\cdot\frac{x^2-x-2}{x^2}=e^{\frac1x}\cdot\frac{(x-2)(x+1)}{x^2}$$
הסימן של
$f'$
הוא הסימן של
$(x-2)(x+1)$,
ולכן:
\begin{itemize}
\item
$f$
עולה ב-$(-\infty,-1]$
וב-$[2,\infty)$.
\item
$f$
יורדת ב-$[-1,0)$
וב-$(0,2]$.
\end{itemize}
$x=-1$
נקודת מקסימום מקומי עם
$f(-1)=\frac1e$,
ו-$x=2$
נקודת מינימום מקומי עם
$f(2)=4\sqrt{e}\approx6.6$.

\textbf{סקיצה:}
משמאל, הגרף צמוד לאסימפטוטה
$y=x+3$,
עולה, חותך את ציר
$x$
ב-$(-2,0)$
עד למקסימום
$\lr{-1,\frac1e}$,
ויורד אל
$0$
כאשר
$x\to0^-$.
מימין, הגרף יורד מ-$\infty$ (ליד ציר
$y$)
עד למינימום
$(2,4\sqrt e)$,
ואז עולה ומתקרב לאסימפטוטה
$y=x+3$.
\end{solbox}
\eMs
\eE
